% paper2_gravity_tau.tex
% The Gravitational Refractive Index as Petz Recovery Fidelity:
% Temporal Asymmetry from Spacetime Geometry
% Paper II in the tau framework series
% Author: Sheng-Kai Huang
% Compile with: pdflatex paper2_gravity_tau.tex

\documentclass[prl,twocolumn,superscriptaddress,longbibliography,floatfix]{revtex4-2}

\usepackage{amsmath}
\usepackage{amssymb}
\usepackage{amsfonts}
\usepackage{amsthm}
\usepackage{bm}
\usepackage{hyperref}
\usepackage{booktabs}
\usepackage{graphicx}

\newtheorem{theorem}{Theorem}

% Convenience macros (same as Paper 1 / Paper 1b)
\newcommand{\kB}{k_{\mathrm{B}}}
\newcommand{\Tr}{\mathrm{Tr}}
\newcommand{\id}{\mathrm{id}}
\newcommand{\Petz}{\mathcal{R}}
\newcommand{\chan}{\mathcal{N}}
\newcommand{\hilb}{\mathcal{H}}
\newcommand{\code}{\mathcal{C}}
\newcommand{\KL}{D}
\newcommand{\ket}[1]{|#1\rangle}
\newcommand{\bra}[1]{\langle#1|}
\newcommand{\braket}[2]{\langle#1|#2\rangle}
\newcommand{\op}[2]{|#1\rangle\langle#2|}
\newcommand{\abs}[1]{\left|#1\right|}
\newcommand{\norm}[1]{\left\|#1\right\|}
\newcommand{\tauasym}{\tau}

% Paper 2 specific macros
\newcommand{\ceff}{c_{\mathrm{eff}}}
\newcommand{\rs}{r_{\mathrm{s}}}
\newcommand{\Sgrav}{\Sigma_{\mathrm{grav}}}
\newcommand{\Fbound}{F_{\mathrm{bound}}}
\newcommand{\nref}{n_{\mathrm{grav}}}

\begin{document}

\title{The Gravitational Refractive Index as Petz Recovery Fidelity:\\
Temporal Asymmetry from Spacetime Geometry}

\author{Sheng-Kai Huang}
\email{akai@fawstudio.com}
\affiliation{Independent Researcher}

\date{\today}

\begin{abstract}
In a companion paper (Paper~I), we established the temporal asymmetry
parameter $\tauasym = 1 - F$ from the Petz recovery framework and the
universal bound $F \geq e^{-\Sigma/2}$.
Here we apply this framework to gravitational fields.
Defining the gravitational entropy production
$\Sgrav \equiv 2|\Phi|/c^2 = \rs/r$, we obtain the identification
$\ceff/c = e^{\Phi/c^2} = e^{-\Sgrav/2} = \Fbound$:
the coordinate speed of light \emph{is} the Petz recovery fidelity
bound [\textsc{new synthesis}].
This connects Dicke's gravitational refractive index to quantum
information recovery.
Three independent derivations---modular flow, gravitational Landauer
principle, and quantum channel analysis---converge on
$\Sgrav = -\ln(-g_{00})$.
Saturation of the Petz bound uniquely determines the exponential
metric $g_{00} = -e^{-\rs/r}$, which has no event horizon
($\tauasym < 1$ everywhere).
We extend the framework to gravitational decoherence (Pikovski
mechanism), the Unruh effect, Bekenstein bound, black hole
complementarity, the Page curve, the Hayden--Preskill protocol,
and ER=EPR, classifying each result as [\textsc{known}],
[\textsc{new synthesis}], [\textsc{new definition}], or
[\textsc{speculative}].
Testable strong-field predictions are tabulated.
\end{abstract}

\maketitle

%=======================================================================
\section{Introduction}
\label{sec:intro}
%=======================================================================

In Paper~I~\cite{Paper1}, we established three results.
First, the Petz recovery map~\cite{Petz1988} is the unique
Bayesian retrodiction functor~\cite{Parzygnat2023}.
Second, the temporal asymmetry parameter
$\tauasym \equiv 1 - F(\rho,\,\Petz_{\sigma,\chan}\!\circ\!\chan(\rho))
\in [0,1]$ quantifies the arrow of time:
$\tauasym = 0$ for closed (unitary) evolution,
$\tauasym > 0$ for open systems, with the universal
bound~\cite{Fawzi2015,JRSWW2018}
\begin{equation}
  \tauasym \;\leq\; 1 - e^{-\Sigma/2}\,,
  \label{eq:tau_bound_paper1}
\end{equation}
where $\Sigma$ is the entropy production of the channel $\chan$.
Third, the equivalence chain
$\tauasym = 0 \Leftrightarrow \Sigma = 0 \Leftrightarrow
I(A;E|B) = 0 \Leftrightarrow \text{QMC}$
connects quantum erasure, error correction, and thermodynamic
irreversibility.

General relativity admits a remarkable reformulation:
gravity acts as a medium with position-dependent refractive
index~\cite{Dicke1957,Einstein1911}.
The coordinate speed of light in a weak gravitational field is
$\ceff/c = e^{\Phi/c^2}$, where $\Phi < 0$ is the Newtonian
potential.
In this paper we show that this expression is not merely a kinematic
accident but a manifestation of the Petz recovery bound:
\begin{equation}
  \frac{\ceff}{c}
  \;=\; e^{\Phi/c^2}
  \;=\; e^{-\Sgrav/2}
  \;=\; \Fbound\,.
  \label{eq:preview}
\end{equation}
Light slowing down in a gravitational field \emph{is} information
loss; the Petz recovery fidelity \emph{is} the refractive index.

%=======================================================================
\section{The Exponential Identification}
\label{sec:identification}
%=======================================================================

\subsection{The coordinate speed of light}

[\textsc{known}]  In the post-Newtonian approximation of GR,
the coordinate speed of light in a static gravitational field
is~\cite{Einstein1911,Dicke1957}
\begin{equation}
  \frac{\ceff}{c} = e^{\Phi/c^2}\,,
  \label{eq:ceff}
\end{equation}
where $\Phi < 0$ is the Newtonian gravitational potential.
This result, first anticipated by Einstein in
1911~\cite{Einstein1911} and placed on rigorous footing by
Dicke~\cite{Dicke1957}, is exact for the exponential
metric and agrees with the Schwarzschild metric through
$O(\rs/r)^2$.

\subsection{Gravitational entropy production}

[\textsc{new definition}]  We define the dimensionless
gravitational entropy production as
\begin{equation}
  \boxed{\;\Sgrav \;\equiv\; \frac{2|\Phi|}{c^2}
  \;=\; \frac{\rs}{r}\,,\;}
  \label{eq:Sgrav_def}
\end{equation}
where $\rs = 2GM/c^2$ is the Schwarzschild radius.
This is an \emph{informational} entropy production---the cost
of retrodiction from infinity to radius $r$---not a
thermodynamic entropy in the Clausius sense.
The Clausius entropy production for a photon in Tolman thermal
equilibrium vanishes identically~\cite{Tolman1930}, because
$T(\Phi)\sqrt{-g_{00}} = T_\infty$.
By contrast, $\Sgrav$ quantifies the informational cost of
climbing the gravitational potential, which persists even in
equilibrium.

Three independent arguments converge on
Eq.~\eqref{eq:Sgrav_def}:

\emph{(a) Modular flow}~[\textsc{known}].---In the
entanglement first-law regime of AQFT on Schwarzschild, the
fractional quantum relative entropy (QRE) loss between radius
$r$ and infinity is $(D_{\mathrm{in}} - D_{\mathrm{out}})/
D_{\mathrm{in}} = \rs/r$, exactly, via the Tolman blueshift
and energy redshift~\cite{DorauMuch2025}.

\emph{(b) Gravitational Landauer principle}~[\textsc{known}].---The
Tolman-modified cost of erasing one bit at radius $r$ exceeds
that at infinity by a factor $1/\sqrt{-g_{00}}$.
The dimensionless excess erasure cost yields
$\Sgrav = -\ln(-g_{00})$~\cite{Herrera2020}, which equals
$\rs/r$ for the exponential metric and
$-\ln(1 - \rs/r) \approx \rs/r + O((\rs/r)^2)$ for
Schwarzschild.

\emph{(c) Quantum channel model}~[\textsc{known}].---Modeling
gravitational redshift as a pure-loss bosonic channel with
intensity transmissivity $\eta = -g_{00}$, the QRE decrease
is $\Sigma = -\ln\eta = \rs/r$~\cite{AhmadiFuentes2014}.

All three routes share the common mathematical core
$\Sgrav = -\ln(-g_{00})$.

\subsection{The triple identification}

[\textsc{new synthesis}]  Substituting
Eq.~\eqref{eq:Sgrav_def} into the Petz fidelity bound
$F \geq e^{-\Sigma/2}$ from Paper~I:
\begin{equation}
  \Fbound = e^{-\Sgrav/2} = e^{\Phi/c^2}\,.
  \label{eq:Fbound}
\end{equation}
Comparing with Eq.~\eqref{eq:ceff}, we obtain the central
result:
\begin{equation}
  \boxed{\;\frac{\ceff}{c}
  \;=\; e^{-\Sgrav/2}
  \;=\; \Fbound\,.\;}
  \label{eq:triple}
\end{equation}
The effective speed of light \emph{is} the Petz recovery
fidelity bound.
Equivalently, the gravitational refractive index is the inverse
fidelity:
\begin{equation}
  \nref(r) = \frac{c}{\ceff} = e^{+\Sgrav/2} = \frac{1}{\Fbound}\,.
  \label{eq:refractive_index}
\end{equation}

Table~\ref{tab:numerical} provides numerical verification
across four orders of magnitude in gravitational compactness.

\begin{table}[t]
\caption{\label{tab:numerical}Gravitational entropy production
$\Sgrav = \rs/r$ and Petz recovery fidelity bound
$\Fbound = e^{-\Sgrav/2}$ for representative astrophysical
objects.  The last row indicates the limit of the weak-field
approximation.}
\begin{ruledtabular}
\begin{tabular}{lccc}
  \textbf{Object} & $\rs/r$ & $\Fbound$ & $\tauasym_{\max}$ \\
  \colrule
  Earth surface    & $1.4{\times}10^{-9}$ & $0.99999999930$
    & $7.0{\times}10^{-10}$ \\
  Sun surface      & $4.2{\times}10^{-6}$ & $0.9999979$
    & $2.1{\times}10^{-6}$ \\
  Neutron star     & $0.4$ & $0.819$
    & $0.181$ \\
  Schwarzschild $r\!=\!\rs$ & $1$ & $0.607$
    & $0.393$ \\
\end{tabular}
\end{ruledtabular}
\end{table}

\emph{Caveat.}---The identification~\eqref{eq:triple} is exact
for the exponential metric
$g_{00} = -e^{-\rs/r}$ and agrees with Schwarzschild through
$O(\rs/r)$.
At $O((\rs/r)^2)$, the two metrics differ:
the Schwarzschild expansion gives
$g_{00}^{\text{Schw}} = -1 + \rs/r - \rs^2/(2r^2) + \cdots$,
identical to the exponential expansion through this order.
The first discrepancy appears at $O((\rs/r)^3)$.
In the solar system ($\rs/r \sim 10^{-6}$), this cubic
correction is of order $10^{-18}$---far below current
measurement capability~\cite{Bertotti2003}.

%=======================================================================
\section{Gravitational Decoherence as Petz Recovery}
\label{sec:decoherence}
%=======================================================================

\subsection{The Pikovski mechanism}

[\textsc{known}]  Pikovski et~al.~\cite{Pikovski2015} showed
that gravitational time dilation induces universal dephasing
of massive objects in spatial superposition.
The Hamiltonian
$H = p^2/(2m) + H_{\mathrm{int}}(1 + \Phi(\hat{x})/c^2)$
produces a position-dependent phase on the internal state,
yielding a CPTP dephasing channel $\chan_{\mathrm{grav}}$ on
the spatial degree of freedom~\cite{Paper1b}.

[\textsc{new synthesis}]  For $N$ internal modes in a thermal
state, the temporal asymmetry parameter is~\cite{Paper1b}
\begin{equation}
  \tauasym_{\mathrm{grav}}(t)
  = 1 - \exp\!\big(-N\Gamma t^2\big)\,,
  \label{eq:tau_pikovski}
\end{equation}
where $\Gamma = (\Delta E)^2(\Delta\Phi)^2/(2\hbar^2 c^4)$
and $\Delta E$ is the internal energy variance.
The connection to Sec.~\ref{sec:identification} is direct:
the $\Sgrav = \rs/r$ from the gravitational potential provides
the entropy production that bounds $\tauasym_{\mathrm{grav}}$
via Eq.~\eqref{eq:tau_bound_paper1}.
The dephasing rate is controlled by
$\Delta\Phi/c^2 \sim \Sgrav/2$, establishing that the
gravitational entropy production governs both the static
(refractive index) and dynamic (decoherence) aspects of
gravitational information loss.

\subsection{Unruh effect}

[\textsc{new synthesis}]  The Unruh effect~\cite{Unruh1976}
states that a uniformly accelerating observer with acceleration
$a$ perceives a thermal bath at temperature
$T_U = \hbar a/(2\pi c \kB)$.
This thermal environment induces entropy production on any
quantum system carried by the observer.
For a system traversing a proper distance $\Delta x$
at acceleration $a$, the dimensionless entropy production
is~\cite{LandiPaternostro2021}
\begin{equation}
  \Sigma_{\mathrm{Unruh}}
  = \frac{2\pi a\,\Delta x}{\hbar c^2}\,,
  \label{eq:Sigma_Unruh}
\end{equation}
giving $\tauasym > 0$ for any accelerating observer.
By the equivalence principle, this connects to gravitational
entropy production: a static observer at radius $r$ experiences
proper acceleration $a = GM/(r^2\sqrt{1 - \rs/r})$, and the
accumulated entropy production over radial displacement
$\Delta r$ reproduces $\Sgrav$ in the appropriate limit.

\subsection{Bekenstein bound in $\tauasym$ language}

[\textsc{new synthesis}]  The Bekenstein
bound~\cite{Bekenstein1981} states that the entropy of a system
of energy $E$ confined to a sphere of radius $R$ satisfies
$S \leq 2\pi ER/(\hbar c)$.
In the $\tauasym$ framework, this becomes
\begin{equation}
  \tauasym \;\leq\; 1 - \exp\!\left(-\frac{\pi E R}{\hbar c}\right),
  \label{eq:bekenstein_tau}
\end{equation}
where we have used $\Sigma_{\mathrm{Bek}} = 2\pi ER/(\hbar c)$
and the Petz bound $\tauasym \leq 1 - e^{-\Sigma/2}$.
This provides a universal upper limit on the temporal asymmetry
of any spatially bounded quantum system.

%=======================================================================
\section{Black Hole Physics through $\tauasym$}
\label{sec:black_holes}
%=======================================================================

\subsection{Complementarity}

[\textsc{new synthesis}]  Black hole complementarity posits that
infalling and external observers describe different but
consistent physics~\cite{Susskind1993}.
In the $\tauasym$ framework, this becomes quantitative:
the infalling observer follows a geodesic and experiences
$\tauasym_{\mathrm{infalling}} \approx 0$ (locally, spacetime
is flat by the equivalence principle), while the external
observer describes the same process with
$\tauasym_{\mathrm{external}} \to 1$ (information appears to
be lost at the horizon).
Complementarity is thus the statement that different quantum
channels---corresponding to different observer
foliations---assign different $\tauasym$ values to the same
physical process.
The Petz recovery map is observer-dependent: the recovery
available to the infalling observer (who has direct access to
the interior) differs from that available to the external
observer (who can only perform operations on Hawking radiation).

\subsection{Page curve as $\tauasym(t)$}

[\textsc{new synthesis}]  The Page curve describes the
entanglement entropy $S_{\mathrm{Page}}(t)$ of Hawking
radiation as a function of time~\cite{Page1993}.
Before the Page time, entropy increases linearly;
after the Page time, it decreases, returning to zero when
evaporation is complete.
In the $\tauasym$ framework, the Page curve translates to
a time-dependent bound on recovery:
\begin{equation}
  \tauasym(t) \;\leq\; 1
  - \exp\!\big(-S_{\mathrm{Page}}(t)/2\big)\,.
  \label{eq:page_tau}
\end{equation}
Before the Page time, $S_{\mathrm{Page}}$ increases and
$\tauasym$ can grow---information becomes progressively harder
to recover.
After the Page time, $S_{\mathrm{Page}}$ decreases and the
bound on $\tauasym$ tightens---information recovery improves.
The Page curve is thus the trajectory of the system through
$\tauasym$ space, and the Page time marks the maximum of
temporal asymmetry.

Recent developments~\cite{Penington2020,AEMM2021} have
provided explicit calculations of the Page curve using the
quantum extremal surface (QES) formula, confirming that
$S_{\mathrm{Page}}(t)$ follows the Page curve and that
unitarity is preserved.
In our framework, the QES formula for generalized entropy
$S_{\mathrm{gen}}$ provides the entropy production
$\Sigma = \Delta S_{\mathrm{gen}}$ that enters the Petz
bound.

\subsection{Hayden--Preskill = Petz recovery}

[\textsc{known}]  Cotler et~al.~\cite{Cotler2019} proved that
the Hayden--Preskill decoder~\cite{HaydenPreskill2007}---the
protocol for recovering information thrown into an old black
hole---\emph{is} the Petz recovery map applied to the
black hole channel.
This is not an analogy: the HP decoder is mathematically
identical to the rotated Petz map
$\Petz_{\sigma,\chan}$~\cite{JRSWW2018}.
The fidelity of information recovery satisfies the same bound
$F \geq e^{-\Sigma/2}$, where $\Sigma$ is the entropy
production of the black hole evaporation channel.

\subsection{Entanglement wedge reconstruction}

[\textsc{known}]  In AdS/CFT, bulk quantum fields are
reconstructed from boundary data using the Petz recovery
map~\cite{ADH2015,CPS2020,Cotler2019}.
The fidelity of entanglement wedge reconstruction
satisfies~\cite{Jafferis2016,Penington2020}
\begin{equation}
  F^2 \;\geq\; \exp(-\Delta S_{\mathrm{gen}})\,,
  \label{eq:ewedge}
\end{equation}
where $\Delta S_{\mathrm{gen}}$ is the change in generalized
entropy.
This is an instance of the same universal recovery
inequality~\cite{JRSWW2018} that underlies our gravitational
identification~\eqref{eq:triple}.

\subsection{ER=EPR}

[\textsc{speculative}]  The ER=EPR
conjecture~\cite{MaldacenaSusskind2013} identifies
Einstein--Rosen bridges (wormholes) with Einstein--Podolsky--Rosen
correlations (entanglement).
In the $\tauasym$ framework, a wormhole connecting two entangled
black holes provides a channel with $\tauasym = 0$ between the
two sides: information can be perfectly recovered through the
bridge, corresponding to the maximal entanglement of the two
black holes.
The non-traversability of the classical ER bridge corresponds
to the fact that the $\tauasym = 0$ channel requires operations
on both sides simultaneously (entanglement does not enable
superluminal signaling).

\subsection{Classification}

Table~\ref{tab:classification} classifies the connections
developed in this section.

\begin{table}[t]
\caption{\label{tab:classification}Classification of results
connecting the $\tauasym$ framework to black hole physics.}
\begin{ruledtabular}
\begin{tabular}{lc}
  \textbf{Connection} & \textbf{Status} \\
  \colrule
  HP decoder = Petz map~\cite{Cotler2019}
    & \textsc{known} \\
  Entanglement wedge via Petz~\cite{CPS2020}
    & \textsc{known} \\
  $F^2 \geq e^{-\Delta S_{\mathrm{gen}}}$~\cite{Penington2020}
    & \textsc{known} \\
  Observer-dependent $\tauasym$ (complementarity)
    & \textsc{new synthesis} \\
  Page curve as $\tauasym(t)$
    & \textsc{new synthesis} \\
  $\ceff/c = \Fbound$ (refractive index)
    & \textsc{new synthesis} \\
  $\Sgrav = \rs/r$ (gravitational entropy production)
    & \textsc{new definition} \\
  ER=EPR as $\tauasym = 0$ channel
    & \textsc{speculative} \\
\end{tabular}
\end{ruledtabular}
\end{table}

%=======================================================================
\section{The Exponential Metric}
\label{sec:exponential}
%=======================================================================

[\textsc{new synthesis}]  If the Petz bound is
\emph{saturated}---i.e., the recovery fidelity equals rather
than exceeds $e^{-\Sgrav/2}$---then the metric component is
determined:
$\sqrt{-g_{00}} = e^{-\Sgrav/2} = e^{-\rs/(2r)}$.
This uniquely determines the exponential
metric~\cite{Papapetrou1954,Yilmaz1958}:
\begin{equation}
  ds^2 = -e^{-\rs/r}\, c^2 dt^2
       + e^{+\rs/r}\!\left(dr^2 + r^2 d\Omega^2\right),
  \label{eq:exp_metric}
\end{equation}
in isotropic coordinates.

\emph{No event horizon.}---Since $e^{-\rs/r} > 0$ for all
$r > 0$, the metric component $g_{00}$ never vanishes.
Equivalently, $\tauasym < 1$ everywhere: complete information
loss never occurs.
This is the information-theoretic statement that event horizons
are forbidden when $\Sgrav < \infty$---a model-independent
consequence of the Petz bound~\cite{JRSWW2018} that does not
depend on the saturation hypothesis.

\emph{Exact solution.}---The exponential metric is not a vacuum
solution of the Einstein equations ($R_{\mu\nu} \neq 0$).
It is an exact solution of Einstein gravity coupled to a phantom
scalar field $\phi = M/r$~\cite{ErnazarovIvashchuk2026}, and
also arises in logarithmic $f(R)$
gravity~\cite{NathSarma2025} and
scalar-Einstein-Gauss-Bonnet theory~\cite{Ernazarov2025}.
Makukov and Mychelkin~\cite{Makukov2020} showed that three
independent scalar-field families within GR all reduce to
the exponential metric when the scalar charge equals the
gravitational mass.
Boonserm et~al.~\cite{Boonserm2018} classified the geometry as
a traversable wormhole.

\emph{Comparison with Schwarzschild.}---The exponential and
Schwarzschild metrics share identical PPN parameters
$\gamma = \beta = 1$ and agree through $O(\rs/r)^2$:
\begin{align}
  g_{00}^{\text{Schw}} &= -1 + \frac{\rs}{r}
    - \frac{\rs^2}{2r^2} + \cdots\,, \\
  g_{00}^{\text{exp}} &= -1 + \frac{\rs}{r}
    - \frac{\rs^2}{2r^2} + \cdots\,.
\end{align}
The first difference appears at $O(\rs/r)^3$:
$-1/6$ (exponential) versus $-1/4$ (Schwarzschild in isotropic
coordinates).
All solar-system tests are passed automatically, including the
Cassini bound $|\gamma - 1| < 2.3\times 10^{-5}$~%
\cite{Bertotti2003}.

\emph{Observational constraints.}---In the strong-field regime,
the exponential metric predicts a photon sphere at $r = \rs$
(isotropic coordinates), a shadow $4.6\%$ larger than
Schwarzschild~\cite{Boonserm2018}, QNM frequencies shifted
by $\sim 4.4\%$, and gravitational-wave echoes at
$\Delta t \approx 4.2\,\rs/c$~\cite{Cardoso2016}.
The EHT measurement of $42 \pm 3\;\mu\text{as}$ for
M87*~\cite{EHT2019} is consistent with both metrics.
NICER X-ray observations~\cite{NICER2019} and next-generation
gravitational-wave detectors can probe the $\sim 5$--$19\%$
strong-field deviations.

\emph{Honest assessment.}---The strong-field regime is
speculative.
The saturation hypothesis ($F = e^{-\Sgrav/2}$ exactly) is
unproven; if nature selects $F > e^{-\Sgrav/2}$, the actual
metric lies between the exponential and Schwarzschild
predictions.
Current data cannot distinguish the two metrics.
The identification~\eqref{eq:triple} is robust in weak fields
and suggestive in strong fields.

%=======================================================================
\section{What $\tauasym$ Can and Cannot Explain}
\label{sec:scope}
%=======================================================================

\emph{Can:}
\begin{itemize}
\item Unify quantum error correction, thermodynamic
  irreversibility, black hole information, and gravitational
  decoherence under a single parameter $\tauasym$.
\item Provide universal bounds: any gravitational channel
  satisfies $F \geq e^{-\Sgrav/2}$, constraining recovery
  independently of the specific protocol.
\item Reproduce the Bekenstein bound, the Unruh entropy
  production, and the Page curve bound in $\tauasym$ language.
\item Identify the Hayden--Preskill decoder as the Petz map
  (established result~\cite{Cotler2019}).
\item Derive the no-horizon theorem: $\tauasym < 1$ for any
  finite-mass object at finite distance
  (model-independent~\cite{JRSWW2018}).
\end{itemize}

\emph{Cannot:}
\begin{itemize}
\item Explain dark matter: at galactic scales,
  $\Sgrav \sim 10^{-6}$, far too small to produce observable
  effects beyond Newtonian gravity.
\item Explain dark energy or the cosmological constant:
  the framework is silent on cosmological dynamics.
\item Derive $g_{ab}$ from $\tauasym$: the relation is
  one-directional ($g \to \Sigma \to \tauasym$); $\tauasym$
  constrains only $g_{00}$, not the full metric tensor.
\item Resolve the firewall paradox~\cite{AMPS2013}: while
  $\tauasym < 1$ removes the event horizon premise, the
  framework does not address the interior structure of compact
  objects.
\item Apply to dynamical spacetimes: the covariant formulation
  $\Fbound(A\!\to\!B) = |\xi|_A/|\xi|_B$ requires a timelike
  Killing vector, which does not exist for gravitational
  collapse or cosmological expansion.
\end{itemize}

%=======================================================================
\section{Discussion and Outlook}
\label{sec:discussion}
%=======================================================================

The identification $\ceff/c = \Fbound$
(Eq.~\eqref{eq:triple}) provides a concrete bridge between
quantum information theory and gravitational physics.
Its logical structure is complementary to the Jacobson
program~\cite{Jacobson1995}: while Jacobson derives
\emph{field equations} from thermodynamic entropy, we derive
the \emph{metric} directly from information recoverability.
Both approaches share the premise that geometry is
information-theoretic in origin; they differ in which
information-theoretic quantity is taken as fundamental and
consequently arrive at different strong-field predictions.

The Dorau--Much result~\cite{DorauMuch2025}---deriving the
semiclassical Einstein equations from quantum relative
entropy---suggests that the two programs may converge:
our $\Sgrav = -\ln(-g_{00})$ and their QRE-based field
equations share the same informational foundation.

Several open questions merit investigation.
(i)~\emph{$\tauasym$ at the Planck scale}: when
$\Sgrav \sim O(1)$, quantum gravity effects become important.
Loop quantum gravity corrections to the exponential metric
have been computed~\cite{GuoYuan2025}, but their
information-theoretic interpretation remains unexplored.
(ii)~\emph{Complexity-weighted $\tauasym$}: the Petz map may be
computationally difficult to implement even when it exists
in principle; a complexity-weighted version of $\tauasym$
could distinguish FAPP from genuine irreversibility.
(iii)~\emph{Non-Markovian gravity}: if the gravitational
channel exhibits memory effects, the entropy production $\Sgrav$
may not be additive, and the Petz bound requires
modification~\cite{LandiPaternostro2021}.
(iv)~\emph{Rotation}: the Kerr analog of the exponential
metric~\cite{Graber2017,MakukovMychelkin2023} has neither
horizons nor ergospheres; its information-theoretic derivation
requires extending the framework beyond spherical symmetry.

In Paper~1b~\cite{Paper1b}, we applied the $\tauasym$ framework
to wavefunction collapse, showing that gravitational
decoherence drives $\tauasym \to 1$ and objectivity ($R \to N$)
at the same timescale.
The present paper extends the framework to strong gravitational
fields.
A future paper (Paper~3) will investigate the weak-field limit
at galactic scales, where quantum corrections to the
gravitational channel may produce observable effects on
rotation curves.

\begin{acknowledgments}
The author thanks the quantum information, quantum gravity,
and quantum foundations communities for the theoretical
foundations on which this work builds.
\end{acknowledgments}

%=======================================================================
% BIBLIOGRAPHY
%=======================================================================

\begin{thebibliography}{35}

% === Paper series ===

\bibitem{Paper1}
S.-K. Huang,
``Universal Recovery: The Petz Map as Retrodiction Functor,
Error Corrector, and Entropy Bound,''
companion paper (2026).

\bibitem{Paper1b}
S.-K. Huang,
``Wavefunction Collapse as Retrodiction Failure: Gravitational
Decoherence, Objectivity, and the $\tau$-$R$ Bridge,''
companion paper (2026).

% === Petz recovery and quantum information ===

\bibitem{Petz1988}
D.~Petz,
``Sufficiency of channels over von Neumann algebras,''
Q.\ J.\ Math.\ \textbf{39}, 97 (1988).

\bibitem{Fawzi2015}
O.~Fawzi and R.~Renner,
``Quantum Conditional Mutual Information and Approximate
Markov Chains,''
Commun.\ Math.\ Phys.\ \textbf{340}, 575 (2015).

\bibitem{JRSWW2018}
M.~Junge, R.~Renner, D.~Sutter, M.~M.~Wilde, and A.~Winter,
``Universal Recovery Maps and Approximate Sufficiency of
Quantum Relative Entropy,''
Ann.\ Henri Poincar\'{e} \textbf{19}, 2955 (2018).

\bibitem{Parzygnat2023}
A.~J. Parzygnat and F.~Buscemi,
``Axioms for retrodiction: Achieving time-reversal symmetry
with a prior,''
Quantum \textbf{7}, 1013 (2023).

% === GR, refractive index, exponential metric ===

\bibitem{Einstein1911}
A.~Einstein,
``\"{U}ber den Einflu{\ss} der Schwerkraft auf die
Ausbreitung des Lichtes,''
Ann.\ Phys.\ (Leipzig) \textbf{35}, 898 (1911).

\bibitem{Dicke1957}
R.~H.~Dicke,
``Gravitation without a Principle of Equivalence,''
Rev.\ Mod.\ Phys.\ \textbf{29}, 363 (1957).

\bibitem{Papapetrou1954}
A.~Papapetrou,
``Eine Theorie des Gravitationsfeldes mit einer Feldfunktion,''
Z.\ Phys.\ \textbf{139}, 518 (1954).

\bibitem{Yilmaz1958}
H.~Yilmaz,
``New Approach to General Relativity,''
Phys.\ Rev.\ \textbf{111}, 1417 (1958).

\bibitem{Makukov2020}
M.~A.~Makukov and E.~G.~Mychelkin,
``Triple Path to the Exponential Metric,''
Found.\ Phys.\ \textbf{50}, 1346 (2020).

\bibitem{Boonserm2018}
P.~Boonserm, T.~Ngampitipan, A.~Simpson, and M.~Visser,
``Exponential metric represents a traversable wormhole,''
Phys.\ Rev.\ D \textbf{98}, 084048 (2018).

\bibitem{Bertotti2003}
B.~Bertotti, L.~Iess, and P.~Tortora,
``A test of general relativity using radio links with the
Cassini spacecraft,''
Nature \textbf{425}, 374 (2003).

\bibitem{ErnazarovIvashchuk2026}
K.~K.~Ernazarov and V.~D.~Ivashchuk,
``Phantom scalar field solution in the exponential metric,''
(2026).

\bibitem{NathSarma2025}
S.~Nath and D.~Sarma,
``Exponential wormhole solution in logarithmic $f(R)$ gravity,''
(2025).

\bibitem{Ernazarov2025}
K.~K.~Ernazarov,
``The Yilmaz-Rosen and Janis-Newman-Winicour metric solutions
in the scalar-Einstein-Gauss-Bonnet 4d gravitational model,''
arXiv:2510.21625 (2025).

\bibitem{Graber2017}
J.~R.~Graber,
``Rotating Yilmaz metric,''
arXiv:1708.05645 (2017).

\bibitem{MakukovMychelkin2023}
M.~A.~Makukov and E.~G.~Mychelkin,
``Rotating exponential metric via Newman--Janis algorithm,''
arXiv:2301.08118 (2023).

% === Gravitational decoherence and thermodynamics ===

\bibitem{Pikovski2015}
I.~Pikovski, M.~Zych, F.~Costa, and \v{C}.~Brukner,
``Universal decoherence due to gravitational time dilation,''
Nat.\ Phys.\ \textbf{11}, 668 (2015).

\bibitem{Tolman1930}
R.~C.~Tolman,
``On the Weight of Heat and Thermal Equilibrium in General
Relativity,''
Phys.\ Rev.\ \textbf{35}, 904 (1930).

\bibitem{LandiPaternostro2021}
G.~T.~Landi and M.~Paternostro,
``Irreversible entropy production: From classical to quantum,''
Rev.\ Mod.\ Phys.\ \textbf{93}, 035008 (2021).

\bibitem{Unruh1976}
W.~G.~Unruh,
``Notes on black-hole evaporation,''
Phys.\ Rev.\ D \textbf{14}, 870 (1976).

\bibitem{Bekenstein1981}
J.~D.~Bekenstein,
``Universal upper bound on the entropy-to-energy ratio for
bounded systems,''
Phys.\ Rev.\ D \textbf{23}, 287 (1981).

% === First-principles derivations of Sigma_grav ===

\bibitem{DorauMuch2025}
L.~Dorau and A.~Much,
``From Quantum Relative Entropy to the Semiclassical Einstein
Equations,''
Phys.\ Rev.\ Lett.\ (2025), arXiv:2510.24491.

\bibitem{Herrera2020}
L.~Herrera,
``Landauer Principle and General Relativity,''
Entropy \textbf{22}, 340 (2020).

\bibitem{AhmadiFuentes2014}
M.~Ahmadi, D.~E.~Bruschi, C.~Sabin, G.~Adesso, and I.~Fuentes,
``Relativistic Quantum Metrology,''
Sci.\ Rep.\ \textbf{4}, 4996 (2014).

% === Black hole information and holography ===

\bibitem{Susskind1993}
L.~Susskind, L.~Thorlacius, and J.~Uglum,
``The Stretched Horizon and Black Hole Complementarity,''
Phys.\ Rev.\ D \textbf{48}, 3743 (1993).

\bibitem{Page1993}
D.~N.~Page,
``Average entropy of a subsystem,''
Phys.\ Rev.\ Lett.\ \textbf{71}, 1291 (1993).

\bibitem{Penington2020}
G.~Penington,
``Entanglement Wedge Reconstruction and the Information Problem,''
J.\ High Energy Phys.\ \textbf{2020}(09), 002 (2020).

\bibitem{AEMM2021}
A.~Almheiri, T.~Hartman, J.~Maldacena, E.~Shaghoulian, and
A.~Tajdini,
``The entropy of Hawking radiation,''
Rev.\ Mod.\ Phys.\ \textbf{93}, 035002 (2021).

\bibitem{HaydenPreskill2007}
P.~Hayden and J.~Preskill,
``Black holes as mirrors: quantum information in random
subsystems,''
J.\ High Energy Phys.\ \textbf{2007}(09), 120 (2007).

\bibitem{Cotler2019}
J.~Cotler, P.~Hayden, G.~Penington, G.~Salton, B.~Swingle,
and M.~Walter,
``Entanglement Wedge Reconstruction via Universal Recovery
Channels,''
Phys.\ Rev.\ X \textbf{9}, 031011 (2019).

\bibitem{ADH2015}
A.~Almheiri, X.~Dong, and D.~Harlow,
``Bulk Locality and Quantum Error Correction in AdS/CFT,''
J.\ High Energy Phys.\ \textbf{2015}(04), 163 (2015).

\bibitem{CPS2020}
H.~Chen, G.~Penington, and G.~Salton,
``Entanglement Wedge Reconstruction using the Petz Map,''
J.\ High Energy Phys.\ \textbf{2020}(01), 168 (2020).

\bibitem{Jafferis2016}
D.~L. Jafferis, A.~Lewkowycz, J.~Maldacena, and S.~S. Suh,
``Relative entropy equals bulk relative entropy,''
J.\ High Energy Phys.\ \textbf{2016}(06), 004 (2016).

\bibitem{MaldacenaSusskind2013}
J.~Maldacena and L.~Susskind,
``Cool horizons for entangled black holes,''
Fortschr.\ Phys.\ \textbf{61}, 781 (2013).

\bibitem{AMPS2013}
A.~Almheiri, D.~Marolf, J.~Polchinski, and J.~Sully,
``Black Holes: Complementarity vs.\ Firewalls,''
J.\ High Energy Phys.\ \textbf{2013}(02), 062 (2013).

% === Gravity from information ===

\bibitem{Jacobson1995}
T.~Jacobson,
``Thermodynamics of Spacetime: The Einstein Equation of State,''
Phys.\ Rev.\ Lett.\ \textbf{75}, 1260 (1995).

% === Observations ===

\bibitem{EHT2019}
Event Horizon Telescope Collaboration,
``First M87 Event Horizon Telescope Results.\ I.,''
Astrophys.\ J.\ Lett.\ \textbf{875}, L1 (2019).

\bibitem{NICER2019}
M.~C.~Miller \emph{et al.},
``PSR J0030+0451 Mass and Radius from NICER Data,''
Astrophys.\ J.\ Lett.\ \textbf{887}, L24 (2019).

\bibitem{Cardoso2016}
V.~Cardoso, E.~Franzato, and P.~Pani,
``Is the Gravitational-Wave Ringdown a Probe of the Event
Horizon?,''
Phys.\ Rev.\ Lett.\ \textbf{116}, 171101 (2016).

\bibitem{DSW2022}
D.~L. Danielson, G.~Satishchandran, and R.~M. Wald,
``Gravitationally mediated entanglement: Newtonian field versus
gravitons,''
Phys.\ Rev.\ D \textbf{105}, 086001 (2022).

% === Recent related work ===

\bibitem{GuoYuan2025}
Y.~Guo and J.~Yuan,
``Quantum effective dynamics of Papapetrou spacetime,''
Eur.\ Phys.\ J.\ C (2025).

\end{thebibliography}

\end{document}
